\section{Discussion}
\label{sec:discussion}

This paper shows that finite binary scan--projection readouts can be folded into stable types through a golden-mean Zeckendorf section.  The central contribution is not merely a reduction in state count.  It is a structured interface: raw microstates, stable normal forms, recursive addresses, visible sigma-algebras, cross-resolution maps, and clarity errors are tracked within one auditable chain.

The mechanism separates stable type formation from raw microstate enumeration.  Zeckendorf legality provides a canonical cross-section for residue classes modulo \(F_{m+2}\).  The folding map projects exponentially many words into stable normal forms, while fold-aware restriction ensures that these normal forms can be compared across adjacent resolutions.  Thus cross-scale stability is not imposed after the fact; it is produced by the interaction between the residue section and the stable inverse-system structure.

The scan error profile \(\varepsilon_m(P;\mu)\) gives a quantitative meaning to finite-resolution clarity.  It is not a heuristic confidence score.  It is the Bayes-optimal binary error among all \(\mathcal F_m\)-measurable decisions.  The cylinder decomposition shows exactly where ambiguity remains: only boundary cylinders contribute to the error.  Boundary-cylinder dimension therefore controls the rate at which deeper scans resolve the event.

\paragraph{Limitations.}
There are three limitations.  First, the golden mean is used here as a canonical normalization grammar and anti-resonant baseline; the paper does not rule out analogous constructions for other Ostrowski, Pisot, or substitution grammars.  Second, the results are mathematical and operational.  They do not constitute an empirical physical theory.  Third, arithmetic, spectral, \v{C}ech-residual, Hilbert-space, quantum, and spacetime interpretations require additional admissibility conditions, limiting procedures, units, and empirical anchoring.  These bridges are outside the main theorem chain.

\paragraph{Future work.}
Future work can proceed in four directions.  The first is an algebraic extension, developing projection-word calculus and stable arithmetic from folded readouts.  The second is a spectral extension, organizing moment-kernel, Hankel rank, primitive orbit, and zeta fingerprints as certificates.  The third is a geometric extension, studying \v{C}ech residuals, holonomy, and discrete curvature candidates.  The fourth is a physicalization audit, investigating whether additional admissibility conditions can turn the mathematical interface into a testable physical model.
